% Generated draft: requires mathematical and source review.
\documentclass[11pt]{article}
\usepackage[margin=25mm]{geometry}
\usepackage{fontspec}
\setmainfont{Latin Modern Roman}
\usepackage{amsmath,amssymb,mathtools,unicode-math}
\setmathfont{Latin Modern Math}
\usepackage{enumitem,needspace,xurl,hyperref}
\setlength{\parindent}{0pt}
\setlength{\parskip}{5pt}
\setlength{\emergencystretch}{3em}
\newcommand{\N}{\mathbb{N}}
\newcommand{\Z}{\mathbb{Z}}
\newcommand{\Q}{\mathbb{Q}}
\newcommand{\R}{\mathbb{R}}
\newcommand{\C}{\mathbb{C}}
\newcommand{\F}{\mathbb{F}}
\newcommand{\A}{\mathbb{A}}
\newcommand{\PP}{\mathbb{P}}
\newcommand{\E}{\mathbb{E}}
\newcommand{\eps}{\varepsilon}
\newcommand{\dd}{\,\mathrm{d}}
\begin{document}
\section*{UnsolvedMath ID: 20002032}
\textbf{Source problem code:} catalogue entry
\subsection{Definitions and mathematical statement}
Problem \(12\) asks what geometric information is encoded by \(Q_3\) or its total integral. The excerpt does not define \(Q_3\). The aid separates a critical/scattering operator from the local Chang--Qing boundary curvature \(T\); they must not be interchanged without specifying the realization and normalization.

For the abstract critical-operator statement, let \(\Sigma^3\) be closed and let \(P_3\) be a regular self-adjoint critical operator with \(P_3 1=0\), the relevant elliptic/Fredholm range and regularity properties, and conformal density law
\[
 q_{\widehat g}=q_g+P_3\omega,\qquad q_g=Q_3(g)\,\dd V_g,
 \quad\widehat g=e^{2\omega}g.
\]
Identify the function kernel conformally as part of the critical setup. Then every moment
\[
 \int_\Sigma u\,q_g\qquad(u\in\ker P_3)
\]
is conformally invariant. Under the stated range hypotheses these moments completely determine the class of \(q_g\) modulo \(\operatorname{im}P_3\): two densities differ by \(P_3\omega\) exactly when all their kernel moments agree. The total \(\int q_g\) alone is complete for this density quotient precisely when \(\ker P_3=\operatorname{span}\{1\}\). This is not a claim that the moments determine the manifold, metric or all filling geometry.

For the Fefferman--Graham scattering realization on the conformal boundary of an admissible four-dimensional conformally compact Einstein filling, the aid reports, in its chosen normalization,
\[
 \int_\Sigma Q_3\,\dd V=3\,V_{\mathrm{ren}},
\]
where \(V_{\mathrm{ren}}\) is the renormalized volume of that chosen filling. The scattering regularity, Einstein normalization and exact construction are not supplied in the record; the formula is retained as a qualified realization-specific assertion.

For the local Chang--Qing realization, let \(g_0\) be the Euclidean metric on the closed unit four-ball and \(r\) its radial coordinate. The aid's test family is
\[
 g_\eps=e^{2\eps(1-r^2)^2}g_0\qquad(\eps\in\R).
\]
The conformal factor and its first normal derivative vanish at \(r=1\), so the induced boundary metric and full second fundamental form stay fixed. Yet the reported total changes are
\[
 \int_{S^3}T_{g_\eps}\,\dd\sigma_{g_\eps}
       -\int_{S^3}T_{g_0}\,\dd\sigma_{g_0}=-48\pi^2\eps,
\]
\[
 \int_{B^4}Q_4(g_\eps)\,\dd V_{g_\eps}
       -\int_{B^4}Q_4(g_0)\,\dd V_{g_0}=96\pi^2\eps.
\]
Thus the bulk change cancels twice the boundary change in the stated convention. Total \(T\) is not determined by boundary metric and second fundamental form alone. These metrics are not asserted to be Einstein fillings, so this example is not substituted into the scattering formula. Full curvature and normal-sign conventions require verification.

The extraction also contains the next participant's background, attributed to Helga Baum. For a spin manifold with spinor bundle \(S\), spin connection \(\nabla^S\), Clifford trace and twistor summand, it describes the Dirac and twistor operators as the two projections of
\[
 \nabla^S:\Gamma(S)\to\Gamma(T^*M\otimes S)
       \cong\Gamma(S)\oplus\Gamma(\mathrm{Tw}).
\]
With the usual normalization, \(D_g\) is Clifford trace and \(P_g\) the complementary twistor projection. Their conformal covariance identifies the respective kernels, so \(h(g)=\dim\ker D_g\) and \(t(g)=\dim\ker P_g\) are invariant for the chosen spin structure and function-space setup. The source says Riemannian cases are better studied than Lorentzian cases. This is adjacent context, not a definition of \(Q_3\); finiteness of a Lorentzian Dirac kernel is not inferred from ellipticity. The page number \(25\) is extraction noise.
\subsection{Short English statement}
The meaning of \(Q_3\) depends on its construction. For a regular self-adjoint critical operator, all kernel moments classify the density modulo conformal gauge; total curvature suffices only for a constant-only kernel. The aid separately relates scattering total curvature to a chosen filling's renormalized volume and gives a local boundary example showing dependence beyond boundary metric and second fundamental form.
\subsection{Sources}
\begin{itemize}
\item Local UnsolvedMath registry, permanent ID \(20002032\), source code \texttt{catalogue entry}; canonical collection \texttt{aim-geometry-notes.json}, notes[369]. Source-item field \(12\), actual Problem \(12\).
\item AIM workshop ``Conformal structure in geometry, analysis, and physics'': \url{https://aimath.org/WWN/confstruct/confstruct.pdf}. The following Helga Baum spinor paragraph is neighboring material and page number \(25\) is not part of the question. No PDF was opened in this session.
\item The aid names Fefferman--Graham scattering and Chang--Qing local boundary curvature, without full authoritative references or complete normalizations. The neighboring text attributes Dirac/twistor discussion to Helga Baum. Original remarks and literature fields are empty; no proposer or year is supplied.
\item Supplied research attempt \(1\), \texttt{partial\_result}, run status \texttt{partially\_solved}, rigor label ``proved,'' novelty confidence ``low.'' Registry classification \texttt{PARTIAL-PROGRESS}, updated 14 August 2026.
\end{itemize}
\end{document}
